% Source: https://artofproblemsolving.com/community/c3h1962331_inscribed_circles
% Stable identity: aops-topic-1962331-q01
We use directed angles between lines, taken modulo $180^\circ$: $\measuredangle(\ell,m)$ is the angle through which line $\ell$ must be turned counterclockwise to become parallel to $m$, and $\measuredangle UVW=\measuredangle(VU,VW)$. With this convention, four points $U,V,W,Z$ (no three collinear) are concyclic if and only if $\measuredangle UZW=\measuredangle UVW$, whatever the order of the points on the circle; and for any triangle $\measuredangle CAB+\measuredangle ABC+\measuredangle BCA=0$.

Let the circumcircles of $ANP$ and $BMP$ meet again at $X$ (if they are tangent at $P$, put $X=P$ and read each line $XP$ below as the common tangent at $P$).

\begin{asy}
usepackage("amsmath");
size(6.5cm);
defaultpen(fontsize(10pt));
pair A=(1.6000,4.2000),B=(0.0000,0.0000),C=(6.0000,0.0000),M=(2.5200,0.0000),N=(4.0200,1.8900),P=(0.9600,2.5200),X=(2.8326,1.1670);
draw(circle((2.6225,2.8486),1.6947),blue);draw(circle((1.2600,0.9629),1.5858),deepgreen);draw(circle((4.2600,0.1593),1.7473),red);
draw(A--B--C--cycle,linewidth(1));
draw(M--X,gray(0.5));draw(N--X,gray(0.5));draw(P--X,gray(0.5));
dot(A);dot(B);dot(C);dot(M);dot(N);dot(P);dot(X);
label("$A$",A,N);label("$B$",B,SW);label("$C$",C,SE);label("$M$",M,S);label("$N$",N,NE);label("$P$",P,W);label("$X$",X,SE);
label("\mbox{}",(min(currentpicture,user=true)+max(currentpicture,user=true))/2,black); // keeps the text after the figure black when LaTeX typesets the labels inline
\end{asy}

Since $X$ lies on circle $ANP$ and $N\in AC$, $P\in AB$,
\[
\measuredangle NXP=\measuredangle NAP=\measuredangle CAB .
\]
Since $X$ lies on circle $BMP$ and $P\in BA$, $M\in BC$,
\[
\measuredangle PXM=\measuredangle PBM=\measuredangle ABC .
\]
Adding,
\[
\measuredangle NXM=\measuredangle NXP+\measuredangle PXM=\measuredangle CAB+\measuredangle ABC=-\measuredangle BCA=\measuredangle ACB=\measuredangle NCM ,
\]
because $N$ lies on line $CA$ and $M$ on line $CB$. Hence $N,X,M,C$ are concyclic, so the circumcircle of $CMN$ also passes through $X$. All three circles pass through $X$ (the \emph{Miquel point}). $\blacksquare$

\textbf{Remark.} The statement is false for \emph{incircles}. For example, if $M,N,P$ are the midpoints of the sides of an equilateral triangle, triangles $ANP$, $BMP$, $CMN$ are the three corner triangles, which meet only at the points $M,N,P$; their incircles lie strictly inside them and so are pairwise disjoint.
